\section{Legal and normative framing}\label{sec:legal}

Brazilian public administration is bound, at every level of the federation, by the five
principles stated in the caput of art.~37 of the 1988 Federal Constitution \citep{brasil1988cf}: legality,
impersonality, morality, publicity and efficiency, the last inserted by Constitutional
Amendment 19/1998.
These are principles rather than rules, and the distinction decides what an event log can be
used for. The distinction is a received one, which we neither restate nor defend: constitutional
doctrine reads principles as optimisation requirements, satisfied in degrees against what is
factually and legally possible, and the term travels in Brazilian scholarship
(\emph{mandados de otimiza\c{c}\~ao}) --- though whether Brazilian administrative doctrine
in fact works in that frame is itself disputed \citep{brgabardo2017eficiencia}, and nothing
below turns on the answer. Only the evidentiary consequence is ours. A rule is
complied with or breached; a principle admits of more and less. Assessing new
instruments of government against administrative-law principles already in force, rather than
against criteria invented for the occasion, has precedent \citep{oswald2018algorithm}. An event
log yields continuous quantities: elapsed
time, idle intervals, handover counts, repetition. Those map onto a gradient, not onto the
per-case legal qualification a breach finding requires, and nothing in this paper rests on a
deadline provision.

Two guarantees in art.~5 add the citizen-facing side. Item LXXVIII, inserted by Amendment
45/2004, assures everyone, in judicial and administrative proceedings alike, reasonable
duration and the means that guarantee celerity of processing. Item XXXIV, ``a'', guarantees
the right of petition to the public authorities in defence of rights or against illegality
or abuse of power, and the caput of that item makes it free of charge. A right of petition
is worth what the petitioner can find out, and knowing that a procedure exists is not
knowing how it runs, how long it takes, or which units it crosses.

At state level, the statute regulating administrative process in the administration studied
here \citep{state2016procedurelaw} sets out in art.~2 a materially broader list of principles
than its federal counterpart, and opens that list with \emph{transparency}, before
\emph{efficiency}, \emph{celerity} and \emph{officiality}, alongside legality, motivation,
proportionality, legal certainty and public interest.
Officiality carries weight here: where the administration must advance a proceeding on its
own initiative, idle time is an administrative fact rather than a party's delay.

Two statutes convert publicity into an operational duty, and it is worth separating what each
one supplies. The access-to-information statute \citep{brasil2011lai} makes publicity the rule
and secrecy the exception, and does not make disclosure wait for a request; its art.~8, \S~3,
III requires public bodies to enable automated access in open, structured and
\emph{machine-readable} formats. That provision governs the \emph{channel and format} of
information a body already holds. It does not by itself say what must be described.

The provision that does is art.~7, \S~2, III of the statute governing users of public services
\citep{brasil2017cdu}, which requires every body to publish a service charter stating, as a
minimum, the \emph{principais etapas para processamento do serviço} --- the main steps by which
a service is processed. The duty is proactive, the online channel is mandated, and the
enumeration is expressly a floor rather than a ceiling. Read with the machine-readability rule
above, the two together amount to a duty to publish an account of how a procedure actually runs,
in a form a machine can consume.

This is the provision the paper turns on, and the reason is a presupposition rather than a
prohibition. A duty to publish the main steps of a procedure presupposes that the administration
knows what those steps are. Neither statute says how that knowledge is to be produced, and
neither supposes it might be absent.

The reach of that duty over this corpus must be stated as an open question rather than assumed,
and it is the sharpest limit on the normative claim. The statute governing users of public
services attaches its duties to \emph{serviços públicos} rendered to a \emph{usuário}, defined
in its art.~2, I as the person who uses or benefits from a public service. The corpus studied
here is predominantly internal administrative procedure whose party is a public servant
(Section~\ref{sec:corpus}). Whether a servant transacting with their own administration is a
\emph{usuário} within that definition is a question of Brazilian administrative law that this
paper does not settle, and on which nothing here should be read as a holding.

There is a serious argument that it is, and we set it out without adopting it. The definition in
art.~2, I is functional rather than status-based: it speaks of the person who \emph{uses or
benefits from} a public service, and does not exclude persons who are also agents of the
administration. A servant applying for leave, contesting a payroll deduction or requesting an
intern is not exercising public authority in that transaction; they are on the receiving end of
a service the administration renders, and the asymmetry of information and of procedural power
that the statute exists to correct is present in the ordinary way. Read functionally, the duty
would reach a substantial part of this corpus. Read as confined to services offered to the
public at large, it would reach only the externally facing minority.

One instrument speaks closer to the point than any commentary, and it is genuinely
double-edged. The normative instrument governing ombudsman units in the federal executive
directs that submissions from internal public agents ``receber\~ao o mesmo tratamento das
manifesta\c{c}\~oes provenientes de usu\'arios de servi\c{c}os p\'ublicos'', and extends to
them ``as mesmas garantias e prote\c{c}\~oes destinadas aos usu\'arios''
\citep{brcgu2024portaria116}. Read one way it tells against the functional reading: a norm that
must \emph{deem} internal agents to be treated as users presupposes that they are not users
already, and an equation enacted for the ombudsman regime does not carry to the publication
duty of its own force. Read the other way it tells for it, since the instrument is made
expressly under the statute, and a regulator has no authority to extend a statute's protections
to a class the statute does not reach at all. We think neither reading is clearly right. What
the provision does establish is that the question is live enough for the federal regulator to
have addressed it.

The same scope limit applies to both of the federal instruments we have found, and it is why
neither closes the question here. The decree regulating the service charter in the federal
executive requires it to state ``as etapas, presentes e futuras, esperadas para a realiza\c{c}\~ao
dos servi\c{c}os, inclu\'idas a estimativas de prazos'' and to provide ``os mecanismos para a
consulta pelos usu\'arios acerca das etapas, cumpridas e pendentes''
\citep{brbrasil2017decreto9094}. That is a closer description of the artefact this paper
produces than anything in the statute itself. But its art.~1 binds \emph{os \'org\~aos e as
entidades do Poder Executivo federal}, and the instrument on ombudsman units is bounded the
same way; against a state secretariat both are persuasive rather than binding. We report the
one that cuts against us and the one that cuts for us on the same footing, because reporting
only the second is what would make a contested question look settled. No decision of the higher
courts or the federal audit body appears to resolve it \citep{brjuruena2021usuario}. We record
the argument and the authority on both sides because a reader should be able to see what turns
on it, and we do not rely on it.

Two further grounds bear on the artefact without depending on how that question is answered, and
their strength must be stated more carefully than the service-charter hook's, because neither
does the same work. The proactive-publicity duty of the access-to-information statute
\citep{brasil2011lai} requires bodies to publish information of collective or general interest
independently of any request. The transparency principle that opens art.~2 of the state statute
governing administrative process \citep{state2016procedurelaw} binds the conduct of
\emph{procedimentos} as such, and is addressed to no class of beneficiary.

Neither, however, generates a duty to \emph{produce} this artefact, and we should not claim they
do. The publicity duty attaches to information a body produces or holds; the record whose
absence is this paper's premise is by hypothesis neither, so the duty reaches the catalogue only
once it exists. And a principle satisfied in degrees cannot compel a particular instrument any
more readily than it can yield the per-case breach finding this section declined at the outset
--- the distinction drawn there binds us here too, and we would be helping ourselves to it
selectively if we let a principle create an obligation it could not create in the other
direction.

What the two grounds do support is narrower and still worth stating. If the catalogue exists, it
is information of the kind the publicity duty covers, so publishing it is not discretionary in
the way publishing an internal working document would be; and the transparency principle gives a
reason, admitting of degrees, to prefer an administration that can account for how its
procedures run over one that cannot. The case for \emph{producing} the artefact is therefore
prudential and normative rather than compelled, except where the service-charter duty applies
and supplies the missing specification of what must be described. We would rather say that than
manufacture an obligation the statutes do not impose.

Data protection constrains the same material from the other side. Art.~12 of the general
data protection statute \citep{brasil2018lgpd} places anonymised data outside the law's
reach unless the anonymisation is reversed, or can be reversed with \emph{esfor\c{c}os
razo\'aveis} (reasonable efforts); arts.~5, III and 5, XI define anonymisation by reference
to technical means that are reasonable and available at the time of processing. That is a
risk standard rather than a binary state, and it makes the strength of an anonymisation
procedure an empirical question about currently available technique. It is the reason this
study treats anonymisation as an experimental arm whose downstream cost is measured, rather
than as a preprocessing detail.

For a reader outside Brazil the nearest conceptual comparator is the right to good
administration in art.~41 of the Charter of Fundamental Rights of the European Union, whose
first paragraph requires that a person's affairs be handled impartially, fairly and within a
reasonable time, and whose second adds the right to be heard, access to the file and the duty
to give reasons. The comparison is conceptual and must be qualified: art.~41 is addressed by
its own terms to ``the institutions, bodies, offices and agencies of the Union''
\citep{eu2012charter}, and the Court of Justice has held that it is accordingly not addressed
to the Member States \citep[para.~67]{cjeu2014ys}, a reading since confirmed
\citep[para.~44]{cjeu2014mukarubega}. The qualification cuts only so far: the same judgments
preserve the right to good administration as a general principle of Union law binding Member
States when they implement it. What does not bind them is art.~41 as such, and neither the
article nor the general principle reaches a non-member federal state. It is offered here to
locate a shared idea --- that handling time and procedural regularity are constitutive of good
administration --- and not as a source of obligation applicable to the administration studied.
The Brazilian arrangement differs in placing efficiency in the constitutional text itself, and,
in the state statute above, in naming celerity and officiality separately from it.

The comparator carries a doctrinal dispute worth importing even though the provision does not
apply. \citet{craig2022article41} reads art.~41 as one right with four limbs, handling within a
reasonable time alongside the hearing, access to the file and the duty to give reasons.
\citet{bousta2013goodadministration} replies that the article, whatever its label, is drafted as
a set of duties on the administration rather than as subjective rights held by individuals. The
dispute need not be settled here, but the two readings license different artefacts, and it is
the second an as-is catalogue fits. Evidence about how procedures actually run bears on whether
an administration is discharging duties of regularity and celerity across its caseload; it
cannot establish that an identified person's right was infringed in a particular case, and this
paper makes no such claim.

Read together, these provisions presuppose something they do not provide: a record of how
administration actually operates. \citet{brauneis2018algorithmic} tested that presupposition
empirically, filing 42 open-records requests across 23 states about government
predictive-algorithm programmes and finding that the recurrent obstacle was not trade
secrecy but that the records had never been created.

The strength of that support should be stated rather than assumed, because it is the paper's
premise and it is carried by a single study from another legal order. \citeauthor{brauneis2018algorithmic}'s
requests concerned government predictive-algorithm programmes --- a novel, procurement-mediated
subject matter --- filed in a request-driven disclosure regime under state statutes with
trade-secret exemptions. The present setting differs on all four counts: routine administrative
procedure, a proactive publication duty, no vendor layer, and a constitutional publicity
principle. The inference that record-absence recurs here is therefore an extrapolation, and we
did not test it: we hold no evidence, and filed no request, establishing that the administration
studied lacked a documented account of its procedures before this work. What we observe is a
weaker and different fact --- that its declared classification concentrates 30\% of cases in one
catch-all (Section~\ref{sec:res-rq1}) --- and a bad classification is not the same failure as a
missing record. Establishing the premise in this jurisdiction would need the empirical
literature on access-to-information compliance here, or a direct enquiry to the body; both are
available and neither is done in this paper. We therefore offer \citeauthor{brauneis2018algorithmic}
as evidence that record-absence is a recurrent failure mode of transparency regimes, not as
evidence about this one.

Our contribution is stated at that level and no higher. An as-is catalogue of process types,
each with a typical path derived from the administration's own event records, supplies one
such record: a statement of the main processing steps in the sense art.~7, \S~2, III
contemplates, machine-readable in the sense art.~8, \S~3, III requires, regenerable from
the log rather than maintained by hand, and inspectable by the staff who operate the
procedures. What it is not should be said with equal precision. The types are descriptive
and revisable, not normative classifications; the catalogue is not evidence of illegality,
not a compliance finding, and not a basis for action against identified staff.

The efficiency principle carries more weight here than a constitutional slogan would, and
Brazilian scholarship states why in terms that bear directly on an event log. The received
account distinguishes efficiency in the conduct of the public agent from efficiency in the
\emph{organisation} of the administration, and it is the second that governs procedure,
because administrative action must be predictable in its course
\citep{brchicoski2004eficiencia}. Applied to administrative process, the same literature
derives obligations that are recognisable in a log: \emph{celeridade} tells against
unnecessary suspensions and a slow tempo, \emph{economia processual} against needless
repetition of acts or phases, and officiality requires that the administration not stand
idle awaiting repeated prompting by the party. Efficiency is treated as a directly effective
constitutional norm rather than ornament, and the controls it licenses are said to include
working methods and procedure itself \citep{brgabardo2017eficiencia, brfranca2000eficiencia}.
Idle intervals, rework loops and stalled cases are therefore not merely operational
defects on which an event log happens to bear; they are the observable form of what this
literature already treats as a legal question. That the scholarship is in Portuguese is a
limitation for readers of this journal, not of the literature.

The federal legislator has since said as much in ordinary law. The statute governing digital
government declares its own object to be ``o aumento da eficiência da administração pública'',
and among its stated principles are transparency in the execution of public services together
with the monitoring of their quality, and --- closest to the artefact proposed here --- ``o uso
da tecnologia para otimizar processos de trabalho da administração pública''
\citep{brasil2021govdigital}. That is a legislative statement that an administration is
expected to use technology on its own working procedures, and not only on the services it
renders outward. Its scope is bounded in the way the other federal instruments discussed above
are: it reaches other federated entities only where they adopt its commands through their own
normative acts, so against the administration studied it is persuasive rather than binding. The
state statute on administrative process, which does bind, opens its list of governing principles
with \emph{transparência} and includes \emph{eficiência}, \emph{celeridade} and
\emph{oficialidade} \citep{state2016procedurelaw}.

Two consequences follow for what this paper is for, and they are worth separating from the
transparency argument. First, if optimising working procedures is a duty rather than a
managerial preference, then a description of how those procedures actually run is not merely a
disclosure artefact: it is the precondition of discharging that duty, since a procedure cannot
be improved before it is known. The typical paths reported here are the object on which such an
analysis would operate --- whether a given path is the one the procedure ought to follow, where
it detours, and which steps could be removed --- and that analysis is future work this paper
enables rather than performs. Second, and more immediately, the electronic system does not
constrain what may happen within a declared process type. It records actions and movements
without enforcing, per type, which sequence is admissible. Nothing in the system prevents two
cases of the same declared type from following entirely different routes. The consequence is
structural rather than incidental: where the procedure is neither constrained by the system nor
described anywhere else, the declared type cannot be relied on as evidence of what was done, and
the only account of the as-is that can be trusted is one derived from the record of what actually
happened. That is the gap this paper addresses, and it is why the catalogue is derived rather
than transcribed from an existing classification.
